% =====================================================================
%  Anatomía de un LLM — Fase I — Hito 1
%  Examen de hito: Fundamentos matemáticos (módulos 1 a 6)
%  Compilar: latexmk -r ../latexmkrc H1-examen.tex
% =====================================================================
\documentclass[11pt,a4paper,twoside,openany]{book}
\usepackage{../estilo-apuntes}

\newcommand{\ApuntesCurso}{Anatomía de un LLM}
\newcommand{\ApuntesModulo}{Hito 1}

\hypersetup{pdftitle={Hito 1. Examen de la Fase I: Fundamentos matemáticos},pdfauthor={Anatomía de un LLM}}

% --- Notación local (la misma que en los apuntes de M03--M06) ----------
\newcommand{\T}{^{\top}}
\newcommand{\KL}[2]{D(#1\,\|\,#2)}
\DeclareMathOperator{\bpb}{bpb}
\DeclareMathOperator{\diag}{diag}
\DeclareMathOperator{\Var}{Var}
\DeclareMathOperator{\rango}{rango}
\DeclareMathOperator{\gen}{gen}
\DeclareMathOperator{\tr}{tr}
\newcommand{\bin}[1]{\textnormal{\texttt{#1}}}

% --- Partes del examen como capítulos con letra ---------------------------
\renewcommand{\thechapter}{\Alph{chapter}}
\newcommand{\parte}[2]{% #1 = título completo, #2 = etiqueta
  \cleardoublepage\refstepcounter{chapter}\label{#2}%
  \chapter*{#1}\addcontentsline{toc}{chapter}{#1}\markboth{#1}{}}

% --- Preguntas del test -------------------------------------------------
\newcounter{pregunta}
\newcommand{\pregunta}[5]{%
  \par\refstepcounter{pregunta}\noindent
  \begin{minipage}{\linewidth}
  \textbf{\thepregunta.}\ #1\enspace{\small(1{,}5 puntos)}
  \begin{enumerate}[label=(\alph*),itemsep=0pt,topsep=2pt,leftmargin=2.4em]
    \item #2
    \item #3
    \item #4
    \item #5
  \end{enumerate}
  \end{minipage}\par\medskip}

% --- Apartados de los problemas y criterios de corrección ----------------
\newcommand{\pts}[1]{\textbf{[#1]}\ }
\newcommand{\apartado}[1]{\par\smallskip\noindent\textbf{Apartado #1}\par\nopagebreak}
\newcommand{\criterio}{\par\noindent\textit{Criterio de corrección.}\ }
\newcommand{\repasar}[1]{\par\smallskip\noindent\textbf{Repasar:} #1}

\begin{document}
\frontmatter
% ------------------------------------------------------------------ Portada
\begin{titlepage}
\centering
\vspace*{2.5cm}
{\large\scshape Anatomía de un LLM}\par
\vspace{0.4cm}
{\normalsize Fundamentos matemáticos de los modelos de lenguaje}\par
\vspace{2.2cm}
\rule{\textwidth}{0.8pt}\par\vspace{0.5cm}
{\small\scshape Fase I \textperiodcentered\ Hito 1 \textperiodcentered\ Examen de hito}\par\vspace{0.3cm}
{\Huge\bfseries Hito 1\par}\vspace{0.3cm}
{\LARGE Fundamentos matemáticos\par}
\vspace{0.3cm}
{\large Módulos 1 a 6\par}
\vspace{0.5cm}\rule{\textwidth}{0.4pt}\par
\vspace{1.6cm}
{\large Examen, proyecto de fase y soluciones}\par
\vfill
{\small Edición de octubre de 2026}\par
\end{titlepage}

% =====================================================================
\chapter*{Instrucciones}
\addcontentsline{toc}{chapter}{Instrucciones}
\markboth{Instrucciones}{}
% =====================================================================

Este examen cierra la Fase~I del curso. Evalúa las herramientas matemáticas que se emplean en el resto del programa: la probabilidad del siguiente token, los vectores y las matrices, los gradientes, la estimación estadística y la medida de la información en bits.

\section*{Condiciones del examen}
\begin{center}
\begin{tabularx}{\textwidth}{@{}lX@{}}
\toprule
Duración & 2 horas. Reparto orientativo: 30 minutos para la parte~A y 90 minutos para la parte~B. \\
Material permitido & Calculadora científica y una hoja A4 de fórmulas escrita a mano por el estudiante. No se permiten apuntes, libros ni dispositivos con conexión. \\
Puntuación total & 100 puntos: parte~A, 30 puntos; parte~B, 70 puntos. \\
Aprobado & 70 puntos o más (70\,\%). \\
Módulos evaluados & M01 a M06 (véase la tabla siguiente). \\
\bottomrule
\end{tabularx}
\end{center}

\section*{Normas}
\begin{enumerate}[label=\arabic*.]
  \item \textbf{Parte A. Cuestiones} (30 puntos): 20 preguntas de opción múltiple, de 1,5 puntos cada una. Cada pregunta tiene una sola respuesta correcta. Las respuestas en blanco y las erróneas valen 0 puntos; las erróneas no restan.
  \item \textbf{Parte B. Problemas} (70 puntos): 6 problemas de desarrollo, uno por módulo. La puntuación de cada apartado figura entre corchetes.
  \item En la parte~B debe escribirse el razonamiento completo. Un resultado correcto sin justificación se puntúa con la mitad de los puntos del apartado.
  \item Los resultados se dan con 3 cifras significativas salvo que se indique otra cosa; en los pasos intermedios conviene conservar 4 o más. Se escribe $\ln$ para el logaritmo natural (resultados en nats) y $\log_2$ para el logaritmo en base~2 (resultados en bits).
  \item El \emph{proyecto de fase} (capítulo final del enunciado) se entrega aparte, en el plazo de una semana, y cuenta en la nota de la fase, no en la de este examen.
\end{enumerate}

\section*{Módulos evaluados}
\begin{center}
\begin{tabular}{@{}cllc@{}}
\toprule
Módulo & Título & Cuestiones & Problema \\
\midrule
M01 & Modelos de lenguaje probabilísticos & 1--4 & B.1 \\
M02 & Álgebra lineal I & 5--6 & B.2 \\
M03 & Álgebra lineal II & 7--11 & B.3 \\
M04 & Cálculo para el aprendizaje automático & 12--14 & B.4 \\
M05 & Probabilidad y estadística & 15--17 & B.5 \\
M06 & Teoría de la información & 18--20 & B.6 \\
\bottomrule
\end{tabular}
\end{center}

\section*{Procedencia de los problemas}
Los problemas B.3, B.4, B.5 y B.6 están adaptados, total o parcialmente, de exámenes y hojas de problemas de cursos universitarios (MIT 18.06, Stanford CS224N, MIT 18.05 y la asignatura \emph{Information Theory} de la Universidad de Cambridge); la fuente se indica al pie de cada enunciado y en la bibliografía. Los enunciados son de redacción propia y los datos se han cambiado para que puedan resolverse a mano. Las soluciones completas están en el \cref{ap:soluciones}, y la tabla de autoevaluación, en el \cref{ap:autoevaluacion}.


\tableofcontents

\mainmatter
% =====================================================================
\parte{Parte A. Cuestiones}{cap:test}
% =====================================================================

Veinte preguntas de opción múltiple, de 1,5 puntos cada una (30 puntos en total). Cada pregunta tiene una sola respuesta correcta; las respuestas erróneas no restan. Tiempo orientativo: 30 minutos.

\subsection*{Módulo 1. Modelos de lenguaje probabilísticos}

\pregunta{La regla de la cadena $P(x_1,\dots,x_T)=\prod_{t=1}^{T}P(x_t\mid x_1,\dots,x_{t-1})$:}
{es una aproximación que solo es válida para modelos de bigramas.}
{es una identidad exacta para cualquier distribución sobre secuencias.}
{exige que los tokens sean independientes entre sí.}
{solo es válida si todas las probabilidades son estrictamente positivas.}

\pregunta{Con un vocabulario de 28 tokens (27 letras y el símbolo \fin), ¿cuántas casillas tiene la tabla de conteos de un modelo de trigramas?}
{$784$}
{$2187$}
{$21\,952$}
{$614\,656$}

\pregunta{Un modelo asigna probabilidad $1/4$ a cada token de un texto de validación. Su perplejidad sobre ese texto es:}
{$0{,}25$}
{$2$}
{$4$}
{$\ln 4\approx 1{,}386$}

\pregunta{Cuando la temperatura tiende a infinito, $T\to\infty$, la distribución $\softmax_T(z)$ de un vector de logits $z\in\R^{K}$ tiende a:}
{un vector one-hot en la posición del logit máximo.}
{la distribución uniforme $(1/K,\dots,1/K)$.}
{$\softmax(z)$, la distribución a temperatura $T=1$.}
{el vector nulo.}

\subsection*{Módulo 2. Álgebra lineal I}

\pregunta{La similitud coseno entre $u=(1,1)$ y $v=(3,0)$ vale:}
{$0$}
{$0{,}5$}
{$0{,}707$}
{$1$}

\pregunta{Sea $E\in\R^{K\times d}$ la matriz de embeddings, con una fila por token, y sea $e_i\in\R^{K}$ el vector one-hot del token $i$. El producto $e_i\T E$ es:}
{la columna $i$ de $E$.}
{la fila $i$ de $E$.}
{la suma de todas las filas de $E$.}
{un escalar.}

\subsection*{Módulo 3. Álgebra lineal II}

\pregunta{En NumPy, un tensor de forma $(8,128,512)$ multiplicado mediante el operador \texttt{@} por una matriz de forma $(512,\,50\,000)$ produce un tensor de forma:}
{$(8,\,128,\,50\,000)$}
{$(8,\,512,\,50\,000)$}
{$(128,\,50\,000)$}
{ninguna: se produce un error de dimensiones.}

\pregunta{La matriz $\begin{pmatrix}2&1\\4&2\end{pmatrix}$:}
{tiene rango~2 y es invertible.}
{tiene rango~1 y uno de sus autovalores es~0.}
{tiene rango~0.}
{tiene determinante~8.}

\pregunta{Una matriz tiene valores singulares $5$, $3$ y $1$. El error, en norma de Frobenius, de su mejor aproximación de rango~1 es:}
{$1$}
{$3$}
{$\sqrt{10}\approx 3{,}16$}
{$4$}

\pregunta{Antes de aplicar el análisis de componentes principales (PCA) se centran los datos porque:}
{sin centrar, la SVD de la matriz de datos no existe.}
{sin centrar, la primera componente tiende a apuntar hacia la media de los datos en lugar de hacia la dirección de máxima varianza.}
{al centrar, los autovalores de la covarianza son negativos.}
{al centrar, la matriz de covarianza es diagonal.}

\pregunta{Se suman dos arrays con las reglas de \emph{broadcasting} de NumPy. ¿Cuál de estas combinaciones de formas produce un error?}
{$(16,128,512)$ y $(512,)$}
{$(16,128,512)$ y $(128,1)$}
{$(16,128,512)$ y $(16,128)$}
{$(16,128,512)$ y $(1,1,512)$}

\subsection*{Módulo 4. Cálculo para el aprendizaje automático}

\pregunta{El gradiente de $f(x,y)=x^{2}y$ en el punto $(1,3)$ es:}
{$(2,1)$}
{$(6,1)$}
{$(6,3)$}
{$(3,2)$}

\pregunta{Sean $p=\softmax(z)$ e $y$ el vector one-hot de la clase correcta $c$. Para la pérdida $L=-\ln\softmax(z)_c$, el gradiente $\nabla_z L$ es:}
{$p$}
{$y-p$}
{$p-y$}
{$-\ln p$}

\pregunta{Para una función $f$ de clase $C^{3}$, el error de truncamiento de la diferencia centrada $[f(x+h)-f(x-h)]/(2h)$ como aproximación de $f'(x)$ es del orden de:}
{$h$}
{$h^{2}$}
{$1/h$}
{no depende de $h$.}

\subsection*{Módulo 5. Probabilidad y estadística}

\pregunta{Una variable de Bernoulli de parámetro $p=0{,}2$ tiene varianza:}
{$0{,}04$}
{$0{,}16$}
{$0{,}2$}
{$0{,}8$}

\pregunta{El estimador de máxima verosimilitud de los parámetros de una distribución categórica a partir de los conteos $N_1,\dots,N_K$ es:}
{$\hat p_j=1/K$}
{$\hat p_j=(N_j+1)/\big(\sum_k N_k+K\big)$}
{$\hat p_j=N_j/\sum_k N_k$}
{$\hat p_j=\sqrt{N_j}/\sum_k\sqrt{N_k}$}

\pregunta{Si se multiplica por~4 el tamaño de un minilote de ejemplos \emph{i.i.d.}, la desviación típica de la pérdida media del minilote:}
{se divide por~4.}
{se divide por~2.}
{no cambia.}
{se multiplica por~2.}

\subsection*{Módulo 6. Teoría de la información}

\pregunta{La entropía de una distribución uniforme sobre 8 símbolos es:}
{1 bit}
{2 bits}
{3 bits}
{8 bits}

\pregunta{Sobre la divergencia de Kullback--Leibler $\KL{p}{q}$ entre distribuciones discretas es cierto que:}
{es simétrica: $\KL{p}{q}=\KL{q}{p}$.}
{puede ser negativa si $q$ se parece mucho a $p$.}
{es mayor o igual que~0 y vale~0 si y solo si $p=q$.}
{es igual a la entropía de $q$.}

\pregunta{Un texto tiene entropía $H(p)=2$ bits por carácter y un modelo $q$ alcanza sobre él una entropía cruzada $H(p,q)=2{,}5$ bits por carácter. La divergencia $\KL{p}{q}$ vale:}
{$0{,}5$ bits}
{$1{,}25$ bits}
{$4{,}5$ bits}
{no puede determinarse con estos datos.}

% =====================================================================
\parte{Parte B. Problemas}{cap:problemas}
% =====================================================================

Seis problemas de desarrollo, uno por módulo (70 puntos en total). La puntuación de cada apartado figura entre corchetes. Tiempo orientativo: 90 minutos.

% ---------------------------------------------------------------------
\begin{ejercicio}{2}{eje:p1}
\textbf{Modelo de bigramas (M01). 12 puntos.}

El corpus de entrenamiento está formado por las palabras «mar», «ama» y «rama». El vocabulario es $\V=\{\fin,\texttt{a},\texttt{m},\texttt{r}\}$, donde \fin\ marca el inicio y el fin de cada palabra.
\begin{enumerate}[label=(\alph*)]
  \item \pts{3} Construya la matriz de conteos $N$ y la matriz de probabilidades de máxima verosimilitud $P$. ¿Cuántos bigramas hay en total?
  \item \pts{3} Calcule $P(\text{mar})$, $P(\text{rama})$ y $P(\text{mama})$. Explique por qué «mama», que no aparece en el corpus, recibe la misma probabilidad que «ama».
  \item \pts{3} Calcule la pérdida media (entropía cruzada) del modelo sobre el corpus de entrenamiento, en nats y en bits, y su perplejidad. Compárela con la del modelo uniforme sobre $\V$ y valore la comparación.
  \item \pts{2} Aplique suavizado de Laplace ($\alpha=1$) y calcule $P(\texttt{a}\mid\texttt{m})$ y $P(\texttt{r}\mid\texttt{r})$.
  \item \pts{1} Con el orden $(\fin,\texttt{a},\texttt{m},\texttt{r})$ y el método de la inversa de la función de distribución, determine el token que sigue a «\texttt{a}» para $u=0{,}35$, $u=0{,}55$ y $u=0{,}9$.
\end{enumerate}
\fuente{elaboración propia, con la notación de Jurafsky y Martin~\cite[cap.~3]{jm3}.}
\end{ejercicio}

% ---------------------------------------------------------------------
\begin{ejercicio}{2}{eje:p2}
\textbf{Embeddings y analogías (M02). 10 puntos.}

Se consideran los siguientes embeddings de juguete en $\R^{3}$; $e_w$ denota el embedding de la palabra $w$.
\begin{center}
\begin{tabular}{@{}lc@{\qquad}lc@{\qquad}lc@{}}
\toprule
Palabra & Vector & Palabra & Vector & Palabra & Vector \\
\midrule
rey & $(3,2,1)$ & mujer & $(1,-2,2)$ & princesa & $(2,-1,2)$ \\
hombre & $(0,2,2)$ & reina & $(3,-2,1)$ & niña & $(0,-1,2)$ \\
\bottomrule
\end{tabular}
\end{center}
\begin{enumerate}[label=(\alph*)]
  \item \pts{2} Calcule $e_{\text{rey}}\cdot e_{\text{reina}}$, $\norm{e_{\text{rey}}}$ y $\cos(e_{\text{rey}},e_{\text{reina}})$.
  \item \pts{4} Calcule $q=e_{\text{rey}}-e_{\text{hombre}}+e_{\text{mujer}}$ y su coseno con $e_{\text{reina}}$, $e_{\text{princesa}}$ y $e_{\text{niña}}$. ¿Qué palabra devuelve la analogía? ¿Qué coordenada parece codificar el género?
  \item \pts{2} Aplique $W=\begin{pmatrix}1&0&0\\0&0&1\end{pmatrix}$ a $e_{\text{rey}}$ y a $e_{\text{reina}}$. ¿Qué coseno tienen las imágenes? Determine el núcleo de $W$ y explique qué información destruye.
  \item \pts{2} \emph{Estimación.} La capa de salida de un LLM multiplica un vector de dimensión $1024$ por una matriz de $50\,000\times 1024$ para obtener los logits. ¿Cuántas operaciones de coma flotante (productos y sumas) cuesta por token? Con un procesador que ejecuta $10^{12}$ operaciones útiles por segundo, ¿cuánto tarda?
\end{enumerate}
\fuente{elaboración propia, a partir del ejemplo de analogías de Mikolov \emph{et al.}~\cite{mikolov-b}.}
\end{ejercicio}

% ---------------------------------------------------------------------
\begin{ejercicio}{2}{eje:p3}
\textbf{Autovalores, SVD y bajo rango (M03). 12 puntos.}

Sean $A=\begin{pmatrix}3&1\\1&3\end{pmatrix}$ y $B=\begin{pmatrix}0&-2\\3&0\end{pmatrix}$.
\begin{enumerate}[label=(\alph*)]
  \item \pts{3} Calcule los autovalores de $A$ y una base ortonormal de autovectores. ¿Son ortogonales los autovectores? Justifique por qué.
  \item \pts{2} Calcule $A^{-1}$, compruebe que $AA^{-1}=I$ y relacione los autovalores de $A^{-1}$ con los de $A$.
  \item \pts{3} Calcule la SVD $B=U\Sigma V\T$ en dos pasos: primero $V$ y $\Sigma$ a partir de $B\T B$; después las columnas de $U$ mediante $u_i=Bv_i/\sigma_i$. Obtenga la mejor aproximación de rango~1, $B_1$, el error $\norm{B-B_1}_F$ y la fracción de energía $\sigma_1^{2}/(\sigma_1^{2}+\sigma_2^{2})$ que retiene. ¿Por qué la SVD de $B$ no puede leerse de sus autovalores, como en el caso de $A$?
  \item \pts{4} \emph{Estimación.} Una matriz de pesos de $2048\times 1024$ se almacena como aproximación de rango $k=64$ (factores $U_k$, $V_k$ y los $k$ valores singulares). ¿Cuántos números se almacenan y qué fracción del original suponen? ¿A partir de qué $k$ deja de compensar? Compárelo con LoRA de rango $r=8$ sobre una matriz de $4096\times 4096$: ¿cuántos parámetros se entrenan y qué porcentaje de la matriz representan?
\end{enumerate}
\fuente{apartado (c) adaptado de MIT 18.06, primavera de 2010, examen~3, problema~3(c)~\cite{mit1806ex3} (SVD en dos pasos a partir de $A\T A$, con otra matriz); apartado (d), del ejercicio de compresión y LoRA de los apuntes del M03 y de Hu \emph{et al.}~\cite{hu-lora}.}
\end{ejercicio}

% ---------------------------------------------------------------------
\begin{ejercicio}{3}{eje:p4}
\textbf{Gradiente de la softmax con entropía cruzada (M04). 14 puntos.}

Una capa lineal sin sesgo produce los logits $z=Wx$, con
\[
W=\begin{pmatrix}1&0\\0&1\\1&-1\end{pmatrix},\qquad x=\begin{pmatrix}1\\2\end{pmatrix},
\]
y el token correcto es el primero ($c=1$). Sean $p=\softmax(z)$, $y$ el vector one-hot de $c$ y $L=-\ln p_c$.
\begin{enumerate}[label=(\alph*)]
  \item \pts{4} \emph{Derivación.} Demuestre que $L$ coincide con la entropía cruzada $H(y,p)=-\sum_k y_k\ln p_k$ y que puede escribirse como $L=-z_c+\ln\sum_j e^{z_j}$. A partir de esta expresión, demuestre que $\partial L/\partial z_k=p_k-y_k$.
  \item \pts{3} Calcule $z$, $p$ y $L$.
  \item \pts{4} Calcule $\nabla_z L$, $\partial L/\partial W$ y $\partial L/\partial x$ mediante la regla de la cadena, indicando la forma de cada resultado.
  \item \pts{2} Dé un paso de descenso de gradiente sobre $W$ con tasa de aprendizaje $\eta=0{,}5$ y calcule la nueva pérdida.
  \item \pts{1} Compruebe $\partial L/\partial W_{11}$ con una diferencia centrada de paso $h=0{,}01$ y calcule el error relativo $r(a,n)=\abs{a-n}/\max(\abs{a},\abs{n})$.
\end{enumerate}
\fuente{apartados (a) y (c) adaptados de Stanford CS224N, tarea~2, preguntas~1(a), 1(b) y 1(d)--(e)~\cite{cs224n-a2} (equivalencia entre la pérdida \emph{naive-softmax} y la entropía cruzada, y gradientes respecto de los vectores y de la matriz de salida); el resto, de elaboración propia.}
\end{ejercicio}

% ---------------------------------------------------------------------
\begin{ejercicio}{2}{eje:p5}
\textbf{Estimación y regla de Bayes (M05). 11 puntos.}
\begin{enumerate}[label=(\alph*)]
  \item \pts{4} \emph{Derivación.} Un modelo muy simple supone que, tras cada letra, el nombre termina con probabilidad constante $q$, de modo que la longitud $L$ (en letras) sigue la distribución geométrica $P(L=k)=(1-q)^{k-1}q$, $k\ge 1$. Se observan las longitudes $3$, $4$, $5$, $5$, $6$ y $7$. Obtenga el estimador de máxima verosimilitud $\hat q$, compruebe que es un máximo y calcule su valor, $P(L=3)$ y $\E[L]=1/\hat q$.
  \item \pts{3} Un cajón contiene tres dados de 4 caras, uno de 6 y uno de 12, todos equilibrados. Se extrae un dado al azar, se lanza una vez y sale un~5. Construya la tabla de Bayes (a priori, verosimilitud, producto y a posteriori) para el tipo de dado. Calcule la probabilidad predictiva a priori de obtener un~5 y la probabilidad de que un segundo lanzamiento del mismo dado dé un~6.
  \item \pts{4} \emph{Estimación numérica.} Un modelo acierta el $78\,\%$ de un banco de pruebas de $600$ preguntas. Calcule el error típico de la exactitud y el intervalo de confianza aproximado del $95\,\%$. ¿Cuántas preguntas harían falta para que la semianchura del intervalo fuera de 1 punto porcentual?
\end{enumerate}
\fuente{apartado (b) adaptado de MIT 18.05, primavera de 2014, examen~2, problema~3~\cite{mit1805ex2} (tabla de Bayes discreta para el tipo de dado, con otros dados y otra pregunta predictiva); apartados (a) y (c), de elaboración propia.}
\end{ejercicio}

% ---------------------------------------------------------------------
\begin{ejercicio}{2}{eje:p6}
\textbf{Entropía, compresión y bits por byte (M06). 11 puntos.}

Una fuente emite los símbolos \texttt{a}, \texttt{b}, \texttt{c}, \texttt{d}, \texttt{e} con probabilidades $p=(1/4,\,1/4,\,1/4,\,1/8,\,1/8)$.
\begin{enumerate}[label=(\alph*)]
  \item \pts{3} Calcule $H(p)$ en bits. Construya un código de Huffman, calcule su longitud media y compárela con $H(p)$. Compruebe la desigualdad de Kraft. Codifique el mensaje «\texttt{abed}» e indique cuántos bits ocupa.
  \item \pts{3} Sean $q$ la distribución uniforme sobre los cinco símbolos y $q'$ la distribución $(1/8,\,1/8,\,1/4,\,1/4,\,1/4)$. Calcule $H(p,q)$, $\KL{p}{q}$, $H(p,q')$ y $\KL{p}{q'}$. ¿Cuántos bits por símbolo de más paga cada modelo? ¿Cuál es mejor?
  \item \pts{5} \emph{Estimación.} Un LLM obtiene una pérdida de validación de $1{,}8$ nats por token, con $4{,}2$ bytes de texto por token en promedio. Calcule su perplejidad, los bits por token y los bits por byte. ¿Cuánto ocuparía 1~GB ($10^{9}$ bytes) de ese texto comprimido mediante codificación aritmética con este modelo? ¿Qué razón de compresión supone?
\end{enumerate}
\fuente{apartado (a) adaptado de la hoja de ejemplos de \emph{Information Theory}, Universidad de Cambridge, curso 2025--26, supervisión~1, ejercicio~12~\cite{camit} (códigos de Huffman y comparación de la longitud media con la entropía), según el planteamiento de Cover y Thomas~\cite[§5.6]{ct}; apartados (b) y (c), de elaboración propia.}
\end{ejercicio}

% =====================================================================
\parte{Proyecto de fase}{cap:proyecto}
% =====================================================================

\section*{Perplejidad: modelo uniforme, bigramas y trigramas}

\paragraph{Objetivo.} Medir, con la métrica con la que se comparan los modelos de lenguaje, cuánto mejora un modelo de caracteres al ampliar el contexto, y explicar por qué más contexto no siempre es mejor cuando hay pocos datos. El proyecto integra el modelo de $n$-gramas y el suavizado (M01, §2.3 y~§3.3), la perplejidad (M01, §5.1), los tensores (M03, §6.1) y la conversión a bits (M06, §2.3).

\paragraph{Datos.} La lista de 316 nombres del cuaderno de prácticas del M01, con el mismo vocabulario de 28 tokens (\fin\ y las 27 letras del alfabeto español) y la misma partición: generador \texttt{np.random.default\_rng(7)}, un $15\,\%$ para validación (47 nombres, 321 tokens que predecir) y los 269 restantes para entrenamiento. No se descarga ningún dato.

\paragraph{Plazo.} Una semana desde la realización del examen. El proyecto cuenta en la nota de la fase, no en la del examen.

\begin{figure}[h]
\centering
\begin{tikzpicture}[node distance=5mm and 7mm,
  caja/.style={draw,rounded corners=2pt,align=center,font=\small,minimum height=8mm,inner sep=3pt},
  flecha/.style={-{Stealth[length=2mm]}}]
\node[caja] (n) {Nombres};
\node[caja,right=of n] (p) {Partición\\entrenamiento / validación};
\node[caja,right=of p] (b) {Bigramas\\$N\in\R^{28\times 28}$};
\node[caja,above=of b] (u) {Uniforme};
\node[caja,below=of b] (t) {Trigramas\\$N_3\in\R^{28\times28\times28}$};
\node[caja,right=of b] (s) {Suavizado $\alpha$\\interpolación $\lambda$};
\node[caja,right=of s] (g) {Pérdida y\\perplejidad\\de validación};
\draw[flecha] (n)--(p);
\draw[flecha] (p)--(b); \draw[flecha] (p.north) |- (u); \draw[flecha] (p.south) |- (t);
\draw[flecha] (b)--(s); \draw[flecha] (t) -| (s);
\draw[flecha] (s)--(g); \draw[flecha] (u) -| (g);
\end{tikzpicture}
\par\smallskip{\small Flujo de trabajo del proyecto de fase.}
\end{figure}

\paragraph{Entregables.}
\begin{enumerate}[label=\arabic*.]
  \item Un cuaderno \texttt{H1-proyecto.ipynb} que se ejecute de principio a fin en CPU en menos de un minuto y sin errores.
  \item Una tabla con la pérdida de validación (en nats y en bits) y la perplejidad de cada modelo.
  \item Una gráfica de la perplejidad de validación frente a $\alpha$, con escala logarítmica en $\alpha$, para bigramas y trigramas.
  \item Un informe breve (a lo sumo una página, en celdas de texto del cuaderno) que responda a las preguntas del paso~6.
\end{enumerate}

\paragraph{Pasos.}
\begin{enumerate}[label=\arabic*.]
  \item \emph{Tokenización con relleno.} Para bigramas, cada nombre se rodea de un símbolo \fin\ a cada lado; para trigramas, se antepone \fin\fin\ y se añade \fin\ al final. Así, ambos modelos predicen exactamente los mismos $\ell+1$ tokens de un nombre de $\ell$ letras. Debe comprobarse con un \texttt{assert} que el número de predicciones coincide.
  \item \emph{Modelo uniforme.} $P(x_t)=1/28$. Justifique, sin hacer ningún cálculo sobre los datos, que su perplejidad es $28$.
  \item \emph{Bigramas.} Con la matriz $N$ y el suavizado de Laplace generalizado del M01, $P_\alpha(j\mid i)=(N_{ij}+\alpha)/\big(\sum_k N_{ik}+28\alpha\big)$.
  \item \emph{Trigramas.} Con el tensor $N_3$ de $28\times 28\times 28$ conteos, $P_\alpha(k\mid i,j)=(N_{3,ijk}+\alpha)/\big(\sum_l N_{3,ijl}+28\alpha\big)$. Cuente cuántas casillas son distintas de cero.
  \item \emph{Barrido de $\alpha$.} Pruebe al menos $\alpha\in\{0{,}01;\ 0{,}05;\ 0{,}1;\ 0{,}2;\ 0{,}5;\ 1\}$ en ambos modelos y elija $\alpha$ en validación. Se recomienda separar una tercera partición de prueba para que la elección de $\alpha$ no contamine la medida final.
  \item \emph{Análisis.} (a) ¿Qué modelo obtiene la menor perplejidad y con qué margen? (b) ¿Por qué el trigrama con suavizado aditivo puede quedar por detrás del bigrama con estos datos? Relaciónelo con la proporción de casillas vacías y con el sobreajuste. (c) ¿Qué perplejidad tendría un modelo perfecto? (d) Convierta la mejor pérdida a bits por carácter y relaciónela con la entropía de la fuente (M06).
  \item \emph{Mejora} (necesaria para la calificación de excelente). Implemente la interpolación $P=\lambda P_3+(1-\lambda)P_2$, con $P_3$ el trigrama sin suavizar y $P_2$ el bigrama suavizado, y elija $\lambda$ en validación (M01, §3.4). Opcionalmente, añada el modelo de unigramas como referencia intermedia.
\end{enumerate}

\paragraph{Valores de referencia.} Con la partición indicada, y solo para comprobar el código: uniforme, $28$; unigramas, $\approx 15{,}7$; bigramas con $\alpha=0{,}2$, $\approx 12{,}4$; trigramas con $\alpha=0{,}2$, $\approx 13{,}2$; interpolación con bigrama de $\alpha=0{,}2$ y $\lambda=0{,}3$, $\approx 11{,}7$. El tensor de trigramas solo tiene 840 casillas no nulas de las $21\,952$.

\section*{Criterios de evaluación}
\begin{center}
\small
\renewcommand{\arraystretch}{1.15}
\begin{tabularx}{\textwidth}{@{}>{\raggedright\arraybackslash}p{2.3cm}c>{\raggedright\arraybackslash}X>{\raggedright\arraybackslash}X>{\raggedright\arraybackslash}X@{}}
\toprule
Criterio & Peso & Excelente & Suficiente & Insuficiente \\
\midrule
Corrección de los modelos & 30\,\% & Los tres modelos reproducen los valores de referencia; hay \texttt{assert} de normalización y del número de predicciones. & Valores correctos, sin comprobaciones explícitas. & Algún modelo mal normalizado, pérdidas infinitas o relleno incoherente. \\
\addlinespace
Metodología de evaluación & 20\,\% & $\alpha$ elegido en validación y medida final en una partición de prueba independiente. & $\alpha$ elegido en validación. & Se evalúa en entrenamiento o se elige $\alpha$ a la vista de la medida final. \\
\addlinespace
Resultados y visualización & 15\,\% & Tabla completa en nats, bits y perplejidad; gráfica con ejes rotulados y escala logarítmica. & Tabla y gráfica legibles, con alguna carencia. & Falta la tabla o la gráfica. \\
\addlinespace
Análisis & 25\,\% & Responde a (a)--(d) con cifras, explica la dispersión del tensor de trigramas e implementa la interpolación. & Responde a (a)--(d) con cifras. & Respuestas vagas o sin números. \\
\addlinespace
Código y reproducibilidad & 10\,\% & Funciones reutilizables, semillas fijas, ejecución en menos de un minuto. & Se ejecuta de principio a fin. & Errores o dependencia del orden manual de ejecución. \\
\bottomrule
\end{tabularx}
\end{center}


\appendix
% =====================================================================
\chapter{Soluciones}
\label{ap:soluciones}
% =====================================================================

Las referencias a los apuntes siguen el formato «M05, §3.2»: módulo, y número de tema y sección tal como aparecen en el PDF de cada módulo. En los pasos intermedios se conservan al menos 4 cifras significativas y los resultados se dan con 3, como piden las instrucciones; cuando existe, se da también el valor exacto como fracción.

% ---------------------------------------------------------------------
\section{Parte A: clave de respuestas}

\begin{center}
\small
\renewcommand{\arraystretch}{1.12}
\begin{tabularx}{\textwidth}{@{}c c >{\raggedright\arraybackslash}X >{\raggedright\arraybackslash}p{2.6cm}@{}}
\toprule
Pregunta & Resp. & Justificación & Apuntes \\
\midrule
1 & (b) & Se obtiene aplicando repetidamente $P(A\cap B)=P(A)\,P(B\mid A)$, sin ninguna aproximación. & M01, §2.2 \\
2 & (c) & Dos tokens de contexto y uno predicho: $28^{3}=21\,952$ casillas. & M01, §2.3 \\
3 & (c) & $\PPL=\exp(-\ln\tfrac14)=4$: el modelo duda por igual entre 4 opciones. & M01, §5.1 \\
4 & (b) & Cada $z_i/T\to 0$, luego cada $e^{z_i/T}\to 1$ y el cociente tiende a $1/K$. & M01, §4.1 \\
5 & (c) & $\cos=3/(\sqrt2\cdot 3)=1/\sqrt2\approx 0{,}707$. & M02, §2.3 \\
6 & (b) & El único~1 de $e_i$ selecciona la fila $i$: es la capa de embeddings. & M02, §5.1 \\
7 & (a) & \texttt{@} actúa sobre las dos últimas dimensiones: $(128,512)(512,50\,000)$ para cada uno de los 8 elementos del lote. & M03, §6.2 \\
8 & (b) & La segunda fila es el doble de la primera: $\det=0$, rango~1 y $\lambda=0$ es autovalor (el otro es la traza, 4). & M03, §1.3 y §3.1 \\
9 & (c) & Eckart--Young: $\big(\sum_{i>1}\sigma_i^{2}\big)^{1/2}=\sqrt{9+1}=\sqrt{10}$. & M03, §5.1 \\
10 & (b) & Sin centrar, $X\T X/n$ es la matriz de segundos momentos y su autovector dominante se orienta hacia la media. & M03, §5.4 \\
11 & (c) & Las formas se alinean por la derecha: $128$ frente a $512$ no coinciden y ninguna es~1. & M03, §6.3 \\
12 & (b) & $\partial_x f=2xy=6$ y $\partial_y f=x^{2}=1$ en $(1,3)$. & M04, §1.2 \\
13 & (c) & Teorema del gradiente de la entropía cruzada; se deriva en el problema B.4(a). & M04, §3.4 \\
14 & (b) & En el desarrollo de Taylor se cancelan los términos pares y queda $\tfrac{h^{2}}{6}f'''(\xi)$. & M04, §6.1 \\
15 & (b) & $\Var(X)=p(1-p)=0{,}2\cdot 0{,}8=0{,}16$. & M05, §2.3 \\
16 & (c) & Maximizar $\sum_j N_j\ln p_j$ con $\sum_j p_j=1$ da contar y normalizar; (b) es el suavizado de Laplace. & M05, §4.2; M01, §3.2 \\
17 & (b) & La desviación típica de la media de $B$ ejemplos es $\sigma/\sqrt B$, y $\sqrt4=2$. & M05, §5.6 \\
18 & (c) & $H=\log_2 8=3$ bits. & M06, §1.2 \\
19 & (c) & Desigualdad de Gibbs; la divergencia no es simétrica. & M06, §2.2 \\
20 & (a) & $H(p,q)=H(p)+\KL{p}{q}$, luego $\KL{p}{q}=2{,}5-2=0{,}5$ bits. & M06, §2.3 \\
\bottomrule
\end{tabularx}
\end{center}

% ---------------------------------------------------------------------
\section{Parte B: soluciones de los problemas}

% =====================================================================  B.1
\begin{solucion}{eje:p1}
\apartado{(a)}
\datos corpus «mar», «ama», «rama»; $\V=\{\fin,\texttt{a},\texttt{m},\texttt{r}\}$.
\planteamiento añadir \fin\ al principio y al final de cada palabra, contar cada par (anterior, siguiente) y dividir cada fila por su total (M01, §3.2).
\pasos
\begin{enumerate}[label=\arabic*.]
  \item Secuencias con marcadores: \fin mar\fin, \fin ama\fin, \fin rama\fin.
  \item Bigramas de \fin mar\fin: $(\fin,\texttt m)$, $(\texttt m,\texttt a)$, $(\texttt a,\texttt r)$, $(\texttt r,\fin)$; son~4.
  \item Bigramas de \fin ama\fin: $(\fin,\texttt a)$, $(\texttt a,\texttt m)$, $(\texttt m,\texttt a)$, $(\texttt a,\fin)$; son~4.
  \item Bigramas de \fin rama\fin: $(\fin,\texttt r)$, $(\texttt r,\texttt a)$, $(\texttt a,\texttt m)$, $(\texttt m,\texttt a)$, $(\texttt a,\fin)$; son~5.
  \item Total: $4+4+5=13$. Conteos y, entre paréntesis, probabilidades (fila: anterior; columna: siguiente):
\end{enumerate}
\begin{center}
\small
\begin{tabular}{@{}lccccc@{}}
\toprule
 & \fin & \texttt a & \texttt m & \texttt r & Total \\
\midrule
\fin & $0$ & $1\ (1/3)$ & $1\ (1/3)$ & $1\ (1/3)$ & $3$ \\
\texttt a & $2\ (0{,}4)$ & $0$ & $2\ (0{,}4)$ & $1\ (0{,}2)$ & $5$ \\
\texttt m & $0$ & $3\ (1)$ & $0$ & $0$ & $3$ \\
\texttt r & $1\ (0{,}5)$ & $1\ (0{,}5)$ & $0$ & $0$ & $2$ \\
\bottomrule
\end{tabular}
\end{center}
\resultado \textbf{13 bigramas; $N$ y $P$ como en la tabla.}
\comprobacion los totales de fila suman $3+5+3+2=13$ y cada fila de $P$ suma~1.
\criterio 2 puntos la matriz $N$; 1 punto la normalización por filas.

\apartado{(b)}
\datos la matriz $P$ del apartado (a).
\planteamiento regla de la cadena con la hipótesis de Markov de orden~1: la probabilidad de una palabra es el producto de $P(\text{siguiente}\mid\text{anterior})$ a lo largo de \fin palabra\fin\ (M01, §2.3).
\pasos
\begin{enumerate}[label=\arabic*.]
  \item $P(\text{mar})=P(\texttt m\mid\fin)P(\texttt a\mid\texttt m)P(\texttt r\mid\texttt a)P(\fin\mid\texttt r)=\tfrac13\cdot1\cdot\tfrac15\cdot\tfrac12=\tfrac1{30}\approx 0{,}03333$.
  \item $P(\text{rama})=\tfrac13\cdot\tfrac12\cdot\tfrac25\cdot1\cdot\tfrac25=\tfrac{2}{75}\approx 0{,}02667$.
  \item $P(\text{mama})=P(\texttt m\mid\fin)P(\texttt a\mid\texttt m)P(\texttt m\mid\texttt a)P(\texttt a\mid\texttt m)P(\fin\mid\texttt a)=\tfrac13\cdot1\cdot\tfrac25\cdot1\cdot\tfrac25=\tfrac{4}{75}\approx 0{,}05333$.
  \item $P(\text{ama})=P(\texttt a\mid\fin)P(\texttt m\mid\texttt a)P(\texttt a\mid\texttt m)P(\fin\mid\texttt a)=\tfrac13\cdot\tfrac25\cdot1\cdot\tfrac25=\tfrac{4}{75}$.
  \item «mama» solo usa transiciones observadas. Como el modelo solo mira el token anterior y $P(\texttt a\mid\texttt m)=1$, el prefijo «ma» cuesta $P(\texttt m\mid\fin)P(\texttt a\mid\texttt m)=\tfrac13$, lo mismo que la «a» inicial de «ama», $P(\texttt a\mid\fin)=\tfrac13$; el resto de ambas palabras es idéntico.
\end{enumerate}
\resultado \textbf{$P(\text{mar})=1/30\approx 0{,}0333$; $P(\text{rama})=2/75\approx 0{,}0267$; $P(\text{mama})=P(\text{ama})=4/75\approx 0{,}0533$.}
\comprobacion $P(\text{mama})/P(\text{ama})=\big(\tfrac13\cdot 1\big)/\tfrac13=1$.
\criterio 1 punto por probabilidad; se resta 1 punto si no se explica la generalización.

\apartado{(c)}
\datos $P(\text{mar})=1/30$, $P(\text{ama})=4/75$, $P(\text{rama})=2/75$; 13 predicciones; $\abs{\V}=4$.
\planteamiento pérdida media $\mathcal J=-\frac1{13}\sum\ln P(\text{palabra})$; en bits, $\mathcal J/\ln 2$; $\PPL=e^{\mathcal J}$ (M01, §5.1).
\pasos
\begin{enumerate}[label=\arabic*.]
  \item $\ln(1/30)=-3{,}4012$; $\ln(4/75)=-2{,}9312$; $\ln(2/75)=-3{,}6243$.
  \item Suma: $-9{,}9567$.
  \item $\mathcal J=9{,}9567/13=0{,}7659$ nats.
  \item En bits: $0{,}7659/0{,}6931=1{,}105$ bits.
  \item $\PPL=e^{0{,}7659}=2{,}151$.
  \item Modelo uniforme: $P=1/4$ en cada paso, $\mathcal J=\ln 4=1{,}386$ nats $=2$ bits y $\PPL=4$.
\end{enumerate}
\resultado \textbf{$\mathcal J\approx 0{,}766$ nats $\approx 1{,}10$ bits; $\PPL\approx 2{,}15$, frente a $4$ del modelo uniforme.} La comparación es optimista: se mide sobre el propio corpus de entrenamiento, no sobre datos de validación (M01, §5.3).
\comprobacion $\PPL=P(\text{corpus})^{-1/13}=\big(\tfrac1{30}\cdot\tfrac4{75}\cdot\tfrac2{75}\big)^{-1/13}=(4{,}741\cdot10^{-5})^{-1/13}=2{,}151$.
\criterio 1 punto la pérdida; 1 punto bits y perplejidad; 1 punto la comparación, con la advertencia sobre el conjunto de entrenamiento.

\apartado{(d)}
\datos fila \texttt m: $(0,3,0,0)$; fila \texttt r: $(1,1,0,0)$; $\abs{\V}=4$; $\alpha=1$.
\planteamiento $P_\alpha(v\mid u)=(N(u,v)+\alpha)/\big(\sum_w N(u,w)+\alpha\abs{\V}\big)$ (M01, §3.3).
\pasos
\begin{enumerate}[label=\arabic*.]
  \item Fila \texttt m suavizada: $(1,4,1,1)$, total $3+4=7$; $P(\texttt a\mid\texttt m)=4/7\approx 0{,}5714$.
  \item Fila \texttt r suavizada: $(2,2,1,1)$, total $2+4=6$; $P(\texttt r\mid\texttt r)=1/6\approx 0{,}1667$.
\end{enumerate}
\resultado \textbf{$P(\bin{a}\mid\bin{m})=4/7\approx 0{,}571$; $P(\bin{r}\mid\bin{r})=1/6\approx 0{,}167$.}
\comprobacion la fila \texttt m suavizada suma $(1+4+1+1)/7=1$; el suavizado reduce la probabilidad máxima ($1\to 0{,}571$) y da probabilidad positiva a los ceros.

\apartado{(e)}
\datos fila \texttt a: $p=(0{,}4;\ 0;\ 0{,}4;\ 0{,}2)$ en el orden $(\fin,\texttt a,\texttt m,\texttt r)$.
\planteamiento se calcula la acumulada $F_k$ y se elige el menor $k$ con $u<F_k$ (M01, §4.3).
\pasos
\begin{enumerate}[label=\arabic*.]
  \item $F=(0{,}4;\ 0{,}4;\ 0{,}8;\ 1)$.
  \item $u=0{,}35<F_1=0{,}4$: sale \fin.
  \item $u=0{,}55$: $F_1=F_2=0{,}4\le 0{,}55<F_3=0{,}8$: sale \texttt m. El token \texttt a, de probabilidad~0, ocupa un intervalo vacío y nunca sale.
  \item $u=0{,}9$: $0{,}8\le 0{,}9<1$: sale \texttt r.
\end{enumerate}
\resultado se obtienen, por este orden, los tokens \fin, \bin{m} y \bin{r}.
\comprobacion los intervalos $[0;\,0{,}4)$, $[0{,}4;\,0{,}8)$ y $[0{,}8;\,1)$ tienen longitudes $0{,}4$, $0{,}4$ y $0{,}2$, iguales a las probabilidades.

\repasar{M01, §2.3 y §3.2 (apartados a y b); M01, §5.1 y §5.3 (c); M01, §3.3 (d); M01, §4.3 (e).}
\end{solucion}

% =====================================================================  B.2
\begin{solucion}{eje:p2}
\apartado{(a)}
\datos $e_{\text{rey}}=(3,2,1)$, $e_{\text{reina}}=(3,-2,1)$.
\planteamiento producto escalar componente a componente, $\norm{v}=\sqrt{v\cdot v}$ y $\cos(u,v)=u\cdot v/(\norm u\norm v)$ (M02, §1.2, §1.3 y §2.3).
\pasos
\begin{enumerate}[label=\arabic*.]
  \item $e_{\text{rey}}\cdot e_{\text{reina}}=9-4+1=6$.
  \item $\norm{e_{\text{rey}}}=\sqrt{9+4+1}=\sqrt{14}\approx 3{,}742$, e igualmente $\norm{e_{\text{reina}}}=\sqrt{14}$.
  \item $\cos=6/14=3/7\approx 0{,}4286$.
\end{enumerate}
\resultado \textbf{producto escalar $6$; $\norm{e_{\text{rey}}}=\sqrt{14}\approx 3{,}74$; $\cos=3/7\approx 0{,}429$.}
\comprobacion $0<3/7<1$: los vectores coinciden en las coordenadas primera y tercera y difieren en el signo de la segunda.

\apartado{(b)}
\datos tabla de embeddings.
\planteamiento aritmética de vectores y coseno con cada candidato (M02, §5.4).
\pasos
\begin{enumerate}[label=\arabic*.]
  \item $q=(3-0+1,\ 2-2-2,\ 1-2+2)=(4,-2,1)$ y $\norm q=\sqrt{21}\approx 4{,}583$.
  \item Cosenos:
\end{enumerate}
\begin{center}
\small
\begin{tabular}{@{}lccc@{}}
\toprule
Palabra $w$ & $q\cdot e_w$ & $\norm{e_w}$ & $\cos(q,e_w)$ \\
\midrule
reina & $12+4+1=17$ & $\sqrt{14}$ & $17/\sqrt{294}\approx 0{,}991$ \\
princesa & $8+2+2=12$ & $3$ & $12/(3\sqrt{21})\approx 0{,}873$ \\
niña & $0+2+2=4$ & $\sqrt5$ & $4/\sqrt{105}\approx 0{,}390$ \\
\bottomrule
\end{tabular}
\end{center}
\begin{enumerate}[label=\arabic*.,start=3]
  \item Las palabras de la consulta tampoco ganan: rey, $9/\sqrt{294}\approx 0{,}525$; mujer, $10/(3\sqrt{21})\approx 0{,}727$; hombre, $-2/(\sqrt{21}\sqrt8)\approx -0{,}154$.
  \item La segunda coordenada vale $+2$ en rey y hombre y $-2$ en reina y mujer: codifica el género.
\end{enumerate}
\resultado \textbf{$q=(4,-2,1)$; la analogía devuelve «reina» ($\cos\approx 0{,}991$); el género está en la segunda coordenada.}
\comprobacion $q-e_{\text{reina}}=(1,0,0)$: $q$ solo difiere de reina en una unidad de la primera coordenada.
\criterio 1 punto $q$; 2 puntos los cosenos; 1 punto la interpretación.

\apartado{(c)}
\datos $W=\begin{pmatrix}1&0&0\\0&0&1\end{pmatrix}$.
\planteamiento $Wv=(v_1,v_3)$: $W$ conserva la primera y la tercera coordenadas (M02, §3.3); el núcleo se define en M03, §1.3.
\pasos
\begin{enumerate}[label=\arabic*.]
  \item $We_{\text{rey}}=(3,1)$ y $We_{\text{reina}}=(3,1)$.
  \item Coseno de las imágenes: $10/10=1$.
  \item $Wv=0\iff v_1=v_3=0$: el núcleo es la recta generada por $(0,1,0)$ y $\rango W=2$.
\end{enumerate}
\resultado \textbf{ambas imágenes son $(3,1)$, con coseno~1. El núcleo de $W$ es $\gen\{(0,1,0)\}$: $W$ destruye la coordenada del género.}
\comprobacion $e_{\text{rey}}-e_{\text{reina}}=(0,4,0)$ está en el núcleo, pues $W(0,4,0)\T=(0,0)\T$.

\apartado{(d)}
\datos vector de dimensión $1024$; matriz de $50\,000\times 1024$; $10^{12}$ operaciones por segundo.
\planteamiento cada logit es un producto escalar de $1024$ términos: $1024$ productos y $1023\approx 1024$ sumas (M02, §4.4).
\pasos
\begin{enumerate}[label=\arabic*.]
  \item Operaciones por logit: $2\cdot 1024=2048$.
  \item Por token: $50\,000\cdot 2048=1{,}024\cdot 10^{8}$.
  \item Tiempo: $1{,}024\cdot10^{8}/10^{12}=1{,}024\cdot 10^{-4}$ s $\approx 0{,}1$ ms.
\end{enumerate}
\resultado \textbf{$\approx 1{,}02\cdot 10^{8}$ operaciones por token, unos $0{,}102$ ms.}
\comprobacion regla de dos operaciones por peso: la matriz tiene $5{,}12\cdot 10^{7}$ pesos y $2\cdot 5{,}12\cdot10^{7}=1{,}024\cdot10^{8}$.
\criterio se acepta cualquier valor entre $5\cdot 10^{7}$ y $1{,}1\cdot 10^{8}$ con razonamiento correcto (contar solo productos da la mitad).

\repasar{M02, §1.2, §1.3 y §2.3 (a); M02, §5.4 (b); M02, §3.3 y M03, §1.3 (c); M02, §4.4 (d).}
\end{solucion}

% =====================================================================  B.3
\begin{solucion}{eje:p3}
\apartado{(a)}
\datos $A=\begin{pmatrix}3&1\\1&3\end{pmatrix}$.
\planteamiento polinomio característico $\det(A-\lambda I)=0$ y resolución de $(A-\lambda I)v=0$ (M03, §3.1).
\pasos
\begin{enumerate}[label=\arabic*.]
  \item $\det(A-\lambda I)=(3-\lambda)^{2}-1=\lambda^{2}-6\lambda+8=(\lambda-4)(\lambda-2)$.
  \item $\lambda_1=4$: $A-4I=\begin{pmatrix}-1&1\\1&-1\end{pmatrix}$ da $v_1=v_2$; vector unitario $q_1=(1,1)/\sqrt2$.
  \item $\lambda_2=2$: $A-2I=\begin{pmatrix}1&1\\1&1\end{pmatrix}$ da $v_1=-v_2$; vector unitario $q_2=(1,-1)/\sqrt2$.
  \item $q_1\cdot q_2=(1-1)/2=0$.
\end{enumerate}
\resultado \textbf{$\lambda=4$ con $(1,1)/\sqrt2$ y $\lambda=2$ con $(1,-1)/\sqrt2$. Son ortogonales porque $A$ es simétrica: por el teorema espectral (M03, §3.3), los autovectores de autovalores distintos de una matriz simétrica son ortogonales.}
\comprobacion $\tr A=6=4+2$ y $\det A=8=4\cdot 2$.
\criterio 1 punto el polinomio y los autovalores; 1 punto los autovectores unitarios; 1 punto la justificación.

\apartado{(b)}
\datos $\det A=8$.
\planteamiento para una matriz $2\times2$, $A^{-1}=\frac1{\det A}\begin{pmatrix}d&-b\\-c&a\end{pmatrix}$ (M03, §2.1).
\pasos
\begin{enumerate}[label=\arabic*.]
  \item $A^{-1}=\frac18\begin{pmatrix}3&-1\\-1&3\end{pmatrix}=\begin{pmatrix}0{,}375&-0{,}125\\-0{,}125&0{,}375\end{pmatrix}$.
  \item $AA^{-1}=\frac18\begin{pmatrix}9-1&-3+3\\3-3&-1+9\end{pmatrix}=I$.
  \item Si $Av=\lambda v$ con $\lambda\neq0$, entonces $A^{-1}v=\lambda^{-1}v$: $A^{-1}$ tiene los mismos autovectores y autovalores $1/4$ y $1/2$.
\end{enumerate}
\resultado \textbf{$A^{-1}=\frac18\begin{pmatrix}3&-1\\-1&3\end{pmatrix}$, con autovalores $1/4$ y $1/2$, inversos de los de $A$.}
\comprobacion $A^{-1}(1,1)\T=\tfrac18(2,2)\T=\tfrac14(1,1)\T$ y $A^{-1}(1,-1)\T=\tfrac18(4,-4)\T=\tfrac12(1,-1)\T$.

\apartado{(c)}
\datos $B=\begin{pmatrix}0&-2\\3&0\end{pmatrix}$.
\planteamiento los $v_i$ y $\sigma_i^{2}$ son los autovectores y autovalores de $B\T B$; las columnas de $U$ se obtienen como $u_i=Bv_i/\sigma_i$ (M03, §4.2). Por Eckart--Young, $B_1=\sigma_1u_1v_1\T$ y $\norm{B-B_1}_F=\sigma_2$ (M03, §5.1).
\pasos
\begin{enumerate}[label=\arabic*.]
  \item $B\T B=\begin{pmatrix}0&3\\-2&0\end{pmatrix}\begin{pmatrix}0&-2\\3&0\end{pmatrix}=\begin{pmatrix}9&0\\0&4\end{pmatrix}$.
  \item Es diagonal: $\sigma_1=\sqrt9=3$ con $v_1=(1,0)$, y $\sigma_2=\sqrt4=2$ con $v_2=(0,1)$; $V=I$ y $\Sigma=\diag(3,2)$.
  \item $u_1=Bv_1/3=(0,3)/3=(0,1)$ y $u_2=Bv_2/2=(-2,0)/2=(-1,0)$; $U=\begin{pmatrix}0&-1\\1&0\end{pmatrix}$, que es ortogonal.
  \item $B_1=3\,u_1v_1\T=3\begin{pmatrix}0&0\\1&0\end{pmatrix}=\begin{pmatrix}0&0\\3&0\end{pmatrix}$.
  \item $B-B_1=\begin{pmatrix}0&-2\\0&0\end{pmatrix}$ y $\norm{B-B_1}_F=2=\sigma_2$.
  \item Energía retenida: $9/(9+4)=9/13\approx 0{,}6923$.
  \item $B$ no es simétrica: su polinomio característico es $\lambda^{2}+6$, con raíces complejas $\pm i\sqrt6$. Sus autovalores no son reales ni sus autovectores ortogonales, de modo que no pueden servir de valores y vectores singulares. Solo para matrices simétricas semidefinidas positivas, como $A$, coinciden la SVD y la descomposición espectral.
\end{enumerate}
\resultado \textbf{$B=\begin{pmatrix}0&-1\\1&0\end{pmatrix}\begin{pmatrix}3&0\\0&2\end{pmatrix}\begin{pmatrix}1&0\\0&1\end{pmatrix}$; $B_1=\begin{pmatrix}0&0\\3&0\end{pmatrix}$; error $2$; energía retenida $9/13\approx 69{,}2\,\%$.}
\comprobacion $U\Sigma V\T=\begin{pmatrix}0\cdot3&-1\cdot2\\1\cdot3&0\cdot2\end{pmatrix}=B$, y $\norm{B}_F^{2}=4+9=13=\sigma_1^{2}+\sigma_2^{2}$.
\criterio 1 punto $V$ y $\Sigma$; 1 punto $U$ y la reconstrucción; 1 punto $B_1$, el error, la energía y la explicación.

\apartado{(d)}
\datos $m\times n=2048\times1024$; $k=64$; LoRA con $r=8$ sobre $d\times d=4096\times4096$.
\planteamiento almacenar el truncamiento de rango $k$ cuesta $k(m+n+1)$ números (M03, §5.2); LoRA entrena $B\in\R^{d_{\text{out}}\times r}$ y $A\in\R^{r\times d_{\text{in}}}$, es decir, $r(d_{\text{in}}+d_{\text{out}})$ parámetros (M03, §5.3).
\pasos
\begin{enumerate}[label=\arabic*.]
  \item Original: $2048\cdot1024=2\,097\,152$ números.
  \item Rango 64: $64\cdot(2048+1024+1)=64\cdot3073=196\,672$ números.
  \item Fracción: $196\,672/2\,097\,152=0{,}09378$, un $9{,}38\,\%$.
  \item Compensa mientras $3073\,k<2\,097\,152$, es decir, $k<682{,}4$: con $k=682$ se almacenan $2\,095\,786$ números y aún se ahorra; con $k=683$, $2\,098\,859$, que ya supera el original.
  \item LoRA: $8\cdot(4096+4096)=65\,536$ parámetros, frente a $4096^{2}=16\,777\,216$.
  \item Porcentaje: $65\,536/16\,777\,216=0{,}003906$, un $0{,}391\,\%$.
\end{enumerate}
\resultado \textbf{$196\,672$ números ($9{,}38\,\%$); compensa hasta $k=682$; LoRA de rango~8 entrena $65\,536$ parámetros, el $0{,}391\,\%$ de la matriz.}
\comprobacion $65\,536/16\,777\,216=2^{16}/2^{24}=2^{-8}=1/256\approx 0{,}39\,\%$.
\criterio 1 punto por cada cifra: almacenamiento, umbral, parámetros de LoRA y porcentaje.

\repasar{M03, §3.1 y §3.3 (a); M03, §2.1 y §3.2 (b); M03, §4.1, §4.2 y §5.1 (c); M03, §5.2 y §5.3 (d).}
\end{solucion}

% =====================================================================  B.4
\begin{solucion}{eje:p4}
\apartado{(a)}
\datos $L=-\ln p_c$, $p=\softmax(z)$, $S=\sum_j e^{z_j}$, $y$ one-hot de $c$.
\planteamiento reescribir la pérdida con la definición de la softmax y derivar término a término (M04, §3.4).
\pasos
\begin{enumerate}[label=\arabic*.]
  \item $H(y,p)=-\sum_k y_k\ln p_k=-\ln p_c$, porque $y_k=0$ para $k\neq c$ e $y_c=1$. Luego $L=H(y,p)$.
  \item $\ln p_c=\ln(e^{z_c}/S)=z_c-\ln S$, de donde $L=-z_c+\ln S$.
  \item $\partial(-z_c)/\partial z_k=-\mathbb I[k=c]=-y_k$.
  \item Regla de la cadena: $\partial(\ln S)/\partial z_k=\frac1S\,\partial S/\partial z_k=e^{z_k}/S=p_k$, pues solo el sumando $j=k$ de $S$ depende de $z_k$.
  \item Sumando: $\partial L/\partial z_k=p_k-y_k$. \qed
\end{enumerate}
\resultado \textbf{$L=H(y,p)=-z_c+\ln\sum_je^{z_j}$ y $\nabla_zL=p-y$.}
\comprobacion $\sum_k(p_k-y_k)=1-1=0$, coherente con la invariancia de la softmax por traslación: sumar una constante a todos los logits no cambia $L$.
\criterio 1 punto la equivalencia con la entropía cruzada; 1 punto la expresión con log-suma-exp; 2 puntos la derivación.

\apartado{(b)}
\datos $W$, $x=(1,2)$, $c=1$.
\planteamiento $z=Wx$, $p=\softmax(z)$, $L=-\ln p_1$.
\pasos
\begin{enumerate}[label=\arabic*.]
  \item $z=(1\cdot1+0\cdot2,\ 0\cdot1+1\cdot2,\ 1\cdot1-1\cdot2)=(1,2,-1)$.
  \item $e^{z}=(2{,}7183;\ 7{,}3891;\ 0{,}36788)$ y $S=10{,}475$.
  \item $p=(0{,}25950;\ 0{,}70538;\ 0{,}035119)$.
  \item $L=-\ln 0{,}25950=1{,}3490$ nats.
\end{enumerate}
\resultado \textbf{$z=(1,2,-1)$; $p\approx(0{,}260;\ 0{,}705;\ 0{,}0351)$; $L\approx 1{,}35$ nats.}
\comprobacion $p$ suma $1{,}0000$, y $L=-z_1+\ln S=-1+2{,}3490=1{,}3490$.

\apartado{(c)}
\datos $p$ del apartado (b), $y=(1,0,0)$, $x=(1,2)$.
\planteamiento $\delta=\nabla_zL=p-y$; para la capa lineal $z=Wx$, $\partial L/\partial W=\delta x\T$ y $\partial L/\partial x=W\T\delta$ (M04, §3.2).
\pasos
\begin{enumerate}[label=\arabic*.]
  \item $\delta=(-0{,}74050;\ 0{,}70538;\ 0{,}035119)$, de forma $3\times1$.
  \item $\partial L/\partial W=\delta x\T$, de forma $(3\times1)(1\times2)=3\times2$, como $W$:
  \[
  \frac{\partial L}{\partial W}=\begin{pmatrix}-0{,}7405&-1{,}4810\\0{,}7054&1{,}4108\\0{,}0351&0{,}0702\end{pmatrix}.
  \]
  \item $\partial L/\partial x=W\T\delta$, de forma $(2\times3)(3\times1)=2\times1$:
  \[
  W\T\delta=\begin{pmatrix}1&0&1\\0&1&-1\end{pmatrix}\begin{pmatrix}-0{,}7405\\0{,}7054\\0{,}0351\end{pmatrix}=\begin{pmatrix}-0{,}7054\\0{,}6703\end{pmatrix}.
  \]
\end{enumerate}
\resultado \textbf{$\nabla_zL\approx(-0{,}741;\ 0{,}705;\ 0{,}0351)$; $\partial L/\partial W$ como arriba ($3\times2$); $\partial L/\partial x\approx(-0{,}705;\ 0{,}670)$ ($2\times1$).}
\comprobacion las componentes de $\delta$ suman $0$, y también cada columna de $\partial L/\partial W$.
\criterio 1 punto $\delta$; 2 puntos $\partial L/\partial W$ con su forma; 1 punto $\partial L/\partial x$.

\apartado{(d)}
\datos $\eta=0{,}5$; $W$ y $\partial L/\partial W$ del apartado (c).
\planteamiento $W'=W-\eta\,\partial L/\partial W$ y nueva evaluación de $z'=W'x$, $p'$ y $L'$ (M04, §1.5).
\pasos
\begin{enumerate}[label=\arabic*.]
  \item $W'=\begin{pmatrix}1+0{,}37025&0+0{,}74050\\0-0{,}35269&1-0{,}70538\\1-0{,}01756&-1-0{,}03512\end{pmatrix}=\begin{pmatrix}1{,}3703&0{,}7405\\-0{,}3527&0{,}2946\\0{,}9824&-1{,}0351\end{pmatrix}$.
  \item $z'_1=1{,}3703+2\cdot0{,}7405=2{,}8513$; $z'_2=-0{,}3527+2\cdot0{,}2946=0{,}2365$; $z'_3=0{,}9824-2\cdot1{,}0351=-1{,}0878$.
  \item $e^{z'}=(17{,}310;\ 1{,}2669;\ 0{,}33696)$, con suma $18{,}913$.
  \item $p'_1=17{,}310/18{,}913=0{,}91520$ y $L'=-\ln0{,}91520=0{,}08861$ nats.
\end{enumerate}
\resultado \textbf{$W'\approx\begin{pmatrix}1{,}37&0{,}741\\-0{,}353&0{,}295\\0{,}982&-1{,}04\end{pmatrix}$; $p'_1\approx0{,}915$; la pérdida baja de $1{,}35$ a $0{,}0886$ nats.}
\comprobacion la aproximación de primer orden predice un descenso $\eta\norm{\partial L/\partial W}_F^{2}=0{,}5\cdot5{,}236=2{,}62$, mayor que el real, $1{,}260$: el paso es grande y la curvatura frena el descenso (M04, §4.4), pero la dirección es de descenso.

\apartado{(e)}
\datos $h=0{,}01$ sobre $W_{11}$; valor analítico $a=p_1-1=-0{,}7405035$.
\planteamiento $n=[L(W_{11}+h)-L(W_{11}-h)]/(2h)$ y error relativo $r(a,n)$ (M04, §6.1 y §6.3). Hacen falta 7 cifras porque se restan números muy próximos.
\pasos
\begin{enumerate}[label=\arabic*.]
  \item $W_{11}=1{,}01$: $z=(1{,}01;\ 2;\ -1)$ y $L_+=1{,}3416168$.
  \item $W_{11}=0{,}99$: $z=(0{,}99;\ 2;\ -1)$ y $L_-=1{,}3564268$.
  \item $n=(1{,}3416168-1{,}3564268)/0{,}02=-0{,}7405020$.
  \item $r=\abs{a-n}/\max(\abs a,\abs n)=1{,}54\cdot10^{-6}/0{,}7405=2{,}08\cdot10^{-6}$.
\end{enumerate}
\resultado \textbf{$n\approx-0{,}740502$ frente a $a\approx-0{,}740504$; $r\approx2{,}1\cdot10^{-6}$: el gradiente analítico es correcto.}
\comprobacion la diferencia es el error de truncamiento previsto, $\tfrac{h^{2}}{6}L'''$, con $L'''=p_1(1-p_1)(1-2p_1)=0{,}09243$ respecto de $W_{11}$: $10^{-4}\cdot0{,}09243/6=1{,}54\cdot10^{-6}$, que coincide con $n-a$. Con $h=0{,}01$ el truncamiento domina sobre el redondeo, por lo que no cabe esperar $r\le10^{-7}$.

\repasar{M04, §3.3 y §3.4 (a); M04, §3.2 (b y c); M04, §1.5 y §4.4 (d); M04, §6.1 y §6.3 (e).}
\end{solucion}

% =====================================================================  B.5
\begin{solucion}{eje:p5}
\apartado{(a)}
\datos $P(L=k)=(1-q)^{k-1}q$; observaciones $3,4,5,5,6,7$, con $n=6$ y $\sum_ik_i=30$.
\planteamiento log-verosimilitud, condición de primer orden y comprobación del signo de la segunda derivada (M05, §4.2).
\pasos
\begin{enumerate}[label=\arabic*.]
  \item $\ell(q)=\sum_i\big[(k_i-1)\ln(1-q)+\ln q\big]=\big(\sum_ik_i-n\big)\ln(1-q)+n\ln q$.
  \item $\ell'(q)=\dfrac nq-\dfrac{\sum_ik_i-n}{1-q}=0\iff n(1-q)=q\big(\sum_ik_i-n\big)\iff\hat q=\dfrac{n}{\sum_ik_i}$.
  \item $\ell''(q)=-\dfrac n{q^{2}}-\dfrac{\sum_ik_i-n}{(1-q)^{2}}<0$ en $(0,1)$: $\ell$ es estrictamente cóncava y $\hat q$ es su máximo global.
  \item $\hat q=6/30=0{,}2$.
  \item $P(L=3)=0{,}8^{2}\cdot0{,}2=0{,}128$ y $\E[L]=1/0{,}2=5$ letras.
\end{enumerate}
\resultado \textbf{$\hat q=n/\sum_ik_i=0{,}2$; $P(L=3)=0{,}128$; $\E[L]=5$ letras.}
\comprobacion la media muestral es $30/6=5=1/\hat q$: el EMV iguala la media teórica a la muestral.
\criterio 2 puntos la derivación (con la comprobación del máximo); 1 punto $\hat q$; 1 punto los dos valores.

\apartado{(b)}
\datos tres dados de 4 caras, uno de 6 y uno de 12; resultado observado: 5.
\planteamiento tabla de Bayes: a posteriori $\propto$ a priori $\times$ verosimilitud; la constante de normalización es la probabilidad predictiva a priori (M05, §1.3).
\pasos
\begin{enumerate}[label=\arabic*.]
  \item Tabla:
\end{enumerate}
\begin{center}
\small
\begin{tabular}{@{}lcccc@{}}
\toprule
Dado & A priori & $P(5\mid\text{dado})$ & Producto & A posteriori \\
\midrule
4 caras & $3/5$ & $0$ & $0$ & $0$ \\
6 caras & $1/5$ & $1/6$ & $1/30$ & $2/3$ \\
12 caras & $1/5$ & $1/12$ & $1/60$ & $1/3$ \\
\midrule
Total & $1$ & & $1/20$ & $1$ \\
\bottomrule
\end{tabular}
\end{center}
\begin{enumerate}[label=\arabic*.,start=2]
  \item Probabilidad predictiva a priori: $P(5)=\tfrac1{30}+\tfrac1{60}=\tfrac1{20}=0{,}05$.
  \item Dado el tipo de dado, los dos lanzamientos son independientes; por probabilidad total con las probabilidades a posteriori, $P(6\text{ en el segundo}\mid5)=\tfrac23\cdot\tfrac16+\tfrac13\cdot\tfrac1{12}=\tfrac4{36}+\tfrac1{36}=\tfrac5{36}\approx0{,}139$.
\end{enumerate}
\resultado \textbf{a posteriori $(0,\ 2/3,\ 1/3)$; $P(5)=1/20$; $P(6\text{ en el segundo lanzamiento}\mid5)=5/36\approx0{,}139$.}
\comprobacion las probabilidades a posteriori suman $2/3+1/3=1$. Antes de observar nada, $P(6)=\tfrac35\cdot0+\tfrac15\cdot\tfrac16+\tfrac15\cdot\tfrac1{12}=\tfrac1{20}=0{,}05$; tras el~5 sube a $5/36\approx0{,}139$, porque la observación descarta los dados de 4 caras, que no pueden dar un~6.
\criterio 1,5 puntos la tabla; 0,5 puntos la predictiva a priori; 1 punto la predictiva a posteriori.

\apartado{(c)}
\datos $\hat p=0{,}78$; $n=600$; $z_{0{,}975}=1{,}96$; semianchura objetivo $0{,}01$.
\planteamiento error típico $\sqrt{\hat p(1-\hat p)/n}$; intervalo $\hat p\pm1{,}96\cdot\mathrm{ET}$ por el teorema central del límite; para el tamaño, despejar $n$ de $1{,}96\sqrt{\hat p(1-\hat p)/n}=0{,}01$ (M05, §5.5).
\pasos
\begin{enumerate}[label=\arabic*.]
  \item $\hat p(1-\hat p)=0{,}78\cdot0{,}22=0{,}1716$.
  \item $\mathrm{ET}=\sqrt{0{,}1716/600}=\sqrt{2{,}86\cdot10^{-4}}=0{,}01691$.
  \item Semianchura: $1{,}96\cdot0{,}01691=0{,}03315$; intervalo $[0{,}7469;\ 0{,}8131]$.
  \item $n=(1{,}96/0{,}01)^{2}\cdot0{,}1716=38\,416\cdot0{,}1716=6592{,}2$, que se redondea hacia arriba: $6593$.
\end{enumerate}
\resultado \textbf{$\mathrm{ET}\approx0{,}0169$; intervalo del $95\,\%$ $\approx[74{,}7\,\%;\ 81{,}3\,\%]$; harían falta $6593$ preguntas.}
\comprobacion reducir la semianchura de $3{,}315$ a $1$ punto exige multiplicar $n$ por $3{,}315^{2}\approx10{,}99$, y $600\cdot10{,}99\approx6592$.
\criterio 1 punto el error típico; 1 punto el intervalo; 2 puntos el tamaño (se aceptan 6592 y 6593).

\repasar{M05, §2.2 y §4.2 (a); M05, §1.3 (b); M05, §5.4 y §5.5 (c).}
\end{solucion}

% =====================================================================  B.6
\begin{solucion}{eje:p6}
\apartado{(a)}
\datos $p=(1/4,1/4,1/4,1/8,1/8)$ para \texttt a, \texttt b, \texttt c, \texttt d, \texttt e; mensaje «\texttt{abed}».
\planteamiento $H(p)=-\sum_ip_i\log_2p_i$ (M06, §1.2); algoritmo de Huffman, que fusiona en cada paso los dos nodos de menor probabilidad (M06, §3.4); desigualdad de Kraft $\sum_i2^{-\ell_i}\le1$ (M06, §3.2).
\pasos
\begin{enumerate}[label=\arabic*.]
  \item $H(p)=3\cdot\tfrac14\cdot2+2\cdot\tfrac18\cdot3=1{,}5+0{,}75=2{,}25$ bits.
  \item Huffman: se fusionan \texttt d y \texttt e ($1/4$); quedan cuatro nodos de $1/4$, que se fusionan por pares, $\{\texttt a,\texttt b\}$ y $\{\texttt c,\texttt{de}\}$ ($1/2$ cada uno), y estos en la raíz.
  \item Código: $\texttt a=\bin{00}$, $\texttt b=\bin{01}$, $\texttt c=\bin{10}$, $\texttt d=\bin{110}$, $\texttt e=\bin{111}$. Cualquier otra elección en los empates da las mismas longitudes $(2,2,2,3,3)$.
  \item Longitud media: $3\cdot\tfrac14\cdot2+2\cdot\tfrac18\cdot3=2{,}25$ bits $=H(p)$.
  \item Kraft: $3\cdot2^{-2}+2\cdot2^{-3}=0{,}75+0{,}25=1$.
  \item «\texttt{abed}» $=\bin{00}\,|\,\bin{01}\,|\,\bin{111}\,|\,\bin{110}=\bin{0001111110}$: $2+2+3+3=10$ bits.
\end{enumerate}
\resultado \textbf{$H(p)=2{,}25$ bits; código $\{\bin{00},\bin{01},\bin{10},\bin{110},\bin{111}\}$ con longitud media $2{,}25$ bits $=H(p)$ y suma de Kraft~1; «\bin{abed}» $=\bin{0001111110}$, 10 bits.}
\comprobacion la longitud media alcanza la entropía porque todas las probabilidades son potencias de $1/2$ y $\ell_i=-\log_2p_i$ (M06, §3.3); además, $-\log_2P(\text{«abed»})=2+2+3+3=10$ bits.
\criterio 1 punto $H(p)$; 1 punto el código y la comprobación de Kraft; 1 punto el mensaje.

\apartado{(b)}
\datos $p$ como en (a); $q=(1/5,\dots,1/5)$; $q'=(1/8,1/8,1/4,1/4,1/4)$.
\planteamiento $H(p,q)=-\sum_ip_i\log_2q_i$ y $\KL{p}{q}=H(p,q)-H(p)$ (M06, §2.2 y §2.3).
\pasos
\begin{enumerate}[label=\arabic*.]
  \item $H(p,q)=\log_25=2{,}3219$ bits y $\KL{p}{q}=2{,}3219-2{,}25=0{,}0719$ bits.
  \item $H(p,q')=\tfrac14\cdot3+\tfrac14\cdot3+\tfrac14\cdot2+\tfrac18\cdot2+\tfrac18\cdot2=2{,}5$ bits y $\KL{p}{q'}=2{,}5-2{,}25=0{,}25$ bits.
\end{enumerate}
\resultado \textbf{$H(p,q)\approx2{,}32$ bits y $\KL{p}{q}\approx0{,}0719$ bits; $H(p,q')=2{,}5$ bits y $\KL{p}{q'}=0{,}25$ bits. Es mejor el modelo uniforme: $q'$ concentra la probabilidad en símbolos equivocados, y un modelo seguro y equivocado puede ser peor que uno que no sabe nada.}
\comprobacion cálculo directo: $\KL{p}{q'}=\tfrac14\log_22+\tfrac14\log_22+\tfrac14\log_21+\tfrac18\log_2\tfrac12+\tfrac18\log_2\tfrac12=0{,}25+0{,}25+0-0{,}125-0{,}125=0{,}25$.
\criterio 0,75 puntos por cada valor.

\apartado{(c)}
\datos $\mathcal J=1{,}8$ nats por token; $\bar b=4{,}2$ bytes por token; $10^{9}$ bytes.
\planteamiento $\PPL=e^{\mathcal J}$; bits por token $=\mathcal J/\ln2$; $\bpb=\mathcal J/(\bar b\ln2)$; tamaño comprimido $=B\cdot\bpb/8$ bytes (M06, §4.3 y §4.4).
\pasos
\begin{enumerate}[label=\arabic*.]
  \item $\PPL=e^{1{,}8}=6{,}050$.
  \item Bits por token: $1{,}8/0{,}69315=2{,}5969$.
  \item Bits por byte: $2{,}5969/4{,}2=0{,}61830$.
  \item Tamaño comprimido: $10^{9}\cdot0{,}61830/8=7{,}729\cdot10^{7}$ bytes $\approx77{,}3$ MB.
  \item Razón de compresión: $8/0{,}61830=12{,}94$.
  \item La codificación aritmética emplea menos de $-\log_2q(x_{1:N})+2$ bits para todo el texto (M06, §4.2): el exceso es despreciable.
\end{enumerate}
\resultado \textbf{$\PPL\approx6{,}05$; $2{,}60$ bits por token; $0{,}618$ bits por byte; 1~GB se comprime a unos $77{,}3$ MB, con una razón de compresión $\approx12{,}9$.}
\comprobacion $10^{9}/12{,}94=7{,}73\cdot10^{7}$ bytes, como en el paso~4.
\criterio 1 punto por cada cifra; se aceptan redondeos razonables.

\repasar{M06, §1.2, §3.2, §3.3 y §3.4 (a); M06, §2.2 y §2.3 (b); M06, §4.2, §4.3 y §4.4, y M01, §5.1 (c).}
\end{solucion}

% =====================================================================
\chapter{Autoevaluación}
\label{ap:autoevaluacion}
% =====================================================================

\section{Cálculo de la nota}

La nota del examen es la suma de los puntos de la parte~A (máximo 30) y de la parte~B (máximo 70). El examen se aprueba con \textbf{70 puntos o más}. Si se obtienen menos de 15 puntos en la parte~A o menos de 40 en la parte~B, conviene repasar esa parte aunque el examen esté aprobado: la Fase~II (módulos 7 a~11) da por sabido todo este material.

\begin{center}
\begin{tabular}{@{}llcc@{}}
\toprule
Parte & Elemento & Puntos & Obtenidos \\
\midrule
A & Cuestiones 1--20 (1,5 cada una) & 30 & \\
\midrule
B & B.1 Modelo de bigramas & 12 & \\
B & B.2 Embeddings y analogías & 10 & \\
B & B.3 Autovalores, SVD y bajo rango & 12 & \\
B & B.4 Gradiente de la softmax con entropía cruzada & 14 & \\
B & B.5 Estimación y regla de Bayes & 11 & \\
B & B.6 Entropía, compresión y bits por byte & 11 & \\
\midrule
 & \textbf{Total} & \textbf{100} & \\
\bottomrule
\end{tabular}
\end{center}

\section{Puntuación por módulo}

La tabla siguiente agrupa los puntos por módulo. Un módulo en el que se obtenga menos del $70\,\%$ de sus puntos debe repasarse aunque el examen esté aprobado.

\begin{center}
\begin{tabular}{@{}lccccc@{}}
\toprule
Módulo & Cuestiones & Puntos A & Problema & Puntos B & Total \\
\midrule
M01 & 1--4 & 6 & B.1 & 12 & 18 \\
M02 & 5--6 & 3 & B.2 & 10 & 13 \\
M03 & 7--11 & 7,5 & B.3 & 12 & 19,5 \\
M04 & 12--14 & 4,5 & B.4 & 14 & 18,5 \\
M05 & 15--17 & 4,5 & B.5 & 11 & 15,5 \\
M06 & 18--20 & 4,5 & B.6 & 11 & 15,5 \\
\midrule
\textbf{Total} & 20 & \textbf{30} & 6 & \textbf{70} & \textbf{100} \\
\bottomrule
\end{tabular}
\end{center}

\section{Puerta de salida de la Fase I}

Para pasar a la Fase~II deben cumplirse las tres condiciones siguientes:
\begin{enumerate}[label=\arabic*.]
  \item un resultado de al menos 9 sobre 12 en el test de cada módulo, M01 a M06;
  \item una autoevaluación del laboratorio interactivo de la fase de al menos el $70\,\%$;
  \item una nota de este examen de hito de al menos el $70\,\%$ (70 puntos sobre 100).
\end{enumerate}
El proyecto de fase se califica con la rúbrica del capítulo correspondiente y cuenta en la nota de la fase.

\section{Qué repasar según los fallos}

\begin{center}
\small
\renewcommand{\arraystretch}{1.1}
\begin{tabularx}{\textwidth}{@{}>{\raggedright\arraybackslash}p{3.6cm}X@{}}
\toprule
Fallo & Secciones de los apuntes \\
\midrule
Cuestiones 1 y 2 & M01, §2.2 (regla de la cadena) y §2.3 ($n$-gramas) \\
Cuestiones 3 y 4 & M01, §5.1 (perplejidad) y §4.1 (softmax con temperatura) \\
B.1 (a), (b) y (e) & M01, §2.3, §3.2 y §4.3 \\
B.1 (c) y (d) & M01, §5.1, §5.3 y §3.3 (suavizado aditivo) \\
Cuestión 5 y B.2 (a) & M02, §1.2, §1.3 y §2.3 \\
Cuestión 6 y B.2 (b) & M02, §5.1 (matriz de embeddings) y §5.4 (analogías) \\
B.2 (c) y (d) & M02, §3.3 y §4.4; M03, §1.3 (núcleo) \\
Cuestiones 7 y 11 & M03, §6.2 (producto por lotes) y §6.3 (\emph{broadcasting}) \\
Cuestión 8 y B.3 (a) y (b) & M03, §1.3, §2.1, §3.1, §3.2 y §3.3 \\
Cuestión 9 y B.3 (c) y (d) & M03, §4.1, §4.2, §5.1, §5.2 y §5.3 \\
Cuestión 10 & M03, §5.4 (PCA) \\
Cuestión 12 & M04, §1.1 y §1.2 \\
Cuestión 13 y B.4 (a)--(c) & M04, §3.2, §3.3 y §3.4 \\
B.4 (d) & M04, §1.5 y §4.4 \\
Cuestión 14 y B.4 (e) & M04, §6.1 y §6.3 \\
Cuestión 15 & M05, §2.3 \\
Cuestión 16 y B.5 (a) & M05, §4.2; M01, §3.2 \\
B.5 (b) & M05, §1.3 \\
Cuestión 17 y B.5 (c) & M05, §5.4, §5.5 y §5.6 \\
Cuestión 18 y B.6 (a) & M06, §1.2, §3.2, §3.3 y §3.4 \\
Cuestiones 19 y 20 y B.6 (b) & M06, §2.2 y §2.3 \\
B.6 (c) & M06, §4.2, §4.3 y §4.4; M01, §5.1 \\
\bottomrule
\end{tabularx}
\end{center}

\section{Qué hacer si no se alcanza el umbral}

\begin{enumerate}[label=\arabic*.]
  \item Anotar cada fallo en el registro de errores, clasificado como error de \emph{concepto}, de \emph{cálculo} o \emph{despiste}.
  \item Para cada módulo por debajo del $70\,\%$, repetir el ciclo de estudio del método del curso con las secciones indicadas en la tabla anterior: teoría, laboratorio, práctica, ejercicios $\star\star\star$ del tema, test y tarjetas de repaso.
  \item Rehacer los problemas fallados sin mirar la solución y, a continuación, los ejercicios del mismo tema en los apuntes del módulo.
  \item Repetir el examen pasada al menos una semana, preferiblemente con datos cambiados, y no avanzar a la Fase~II hasta superar las tres condiciones de la puerta de salida.
\end{enumerate}


\backmatter
% =====================================================================
\begin{thebibliography}{99}
\addcontentsline{toc}{chapter}{Bibliografía}
\markboth{Bibliografía}{}
\small

\bibitem{camit}
R.~Harle.
\newblock Information Theory.
\newblock University of Cambridge, Department of Computer Science and Technology, Part~II, curso 2025--26.
\newblock \emph{Examples Sheet}, supervisión~1, ejercicio~12 (códigos de Huffman): \url{https://www.cl.cam.ac.uk/teaching/2526/InfoTheory/InfoTheoryExamplesSheet.pdf}

\bibitem{ct}
T.~M. Cover y J.~A. Thomas.
\newblock \emph{Elements of Information Theory}, 2.ª ed.
\newblock Wiley, 2006.
\newblock Cap.~5, «Data Compression».

\bibitem{cs224n-a2}
Stanford University.
\newblock CS224N: Natural Language Processing with Deep Learning. Tarea~2, sección~1, «Understanding word2vec», preguntas~1(a), 1(b) y 1(d)--(e).
\newblock Consultada en octubre de 2026.
\newblock \url{https://web.stanford.edu/class/cs224n/assignments/a2.pdf}

\bibitem{hu-lora}
E.~J. Hu, Y.~Shen, P.~Wallis, Z.~Allen-Zhu, Y.~Li, S.~Wang, L.~Wang y W.~Chen.
\newblock LoRA: Low-rank adaptation of large language models.
\newblock En \emph{International Conference on Learning Representations (ICLR)}, 2022.
\newblock \url{https://arxiv.org/abs/2106.09685}

\bibitem{jm3}
D.~Jurafsky y J.~H. Martin.
\newblock \emph{Speech and Language Processing}, 3.ª ed., borrador del 19 de agosto de 2026.
\newblock Cap.~3, «N-gram Language Models».
\newblock \url{https://web.stanford.edu/~jurafsky/slp3/}

\bibitem{mikolov-b}
T.~Mikolov, W.-t. Yih y G.~Zweig.
\newblock Linguistic regularities in continuous space word representations.
\newblock En \emph{Proceedings of NAACL-HLT}, págs.~746--751, 2013.
\newblock \url{https://aclanthology.org/N13-1090/}

\bibitem{mit1805ex2}
J.~Orloff y J.~Bloom.
\newblock 18.05: Introduction to Probability and Statistics.
\newblock Massachusetts Institute of Technology, MIT OpenCourseWare, primavera de 2014. Examen~2, problema~3.
\newblock \url{https://ocw.mit.edu/courses/18-05-introduction-to-probability-and-statistics-spring-2014/pages/exams/}

\bibitem{mit1806ex3}
G.~Strang.
\newblock 18.06: Linear Algebra.
\newblock Massachusetts Institute of Technology, MIT OpenCourseWare, primavera de 2010. Examen~3, problema~3.
\newblock \url{https://ocw.mit.edu/courses/18-06-linear-algebra-spring-2010/pages/exams/}

\end{thebibliography}


\end{document}
